\documentclass[11pt]{article}
\usepackage[a4paper,margin=25mm]{geometry}
\usepackage{amsmath,amssymb,mathtools,bm,hyperref}
\usepackage[T1]{fontenc}
\setlength{\emergencystretch}{2em}
\newcommand{\tr}{\operatorname{tr}}
\newcommand{\Id}{\mathbb{1}}
\newcommand{\LL}{\mathcal{L}}
\newcommand{\PP}{\mathcal{P}}
\newcommand{\GG}{\mathcal{G}}
\newcommand{\ket}[1]{\lvert #1\rangle}
\newcommand{\bra}[1]{\langle #1\rvert}
\newcommand{\comm}[2]{[#1,#2]}
\title{Momentum-resolved tangent excitations of RCMPS\\
Connected environments and a reference implementation}
\author{RcmpsCore refactor implementation notes}
\date{September 2026}
\begin{document}
\maketitle

\begin{abstract}
We derive an implementation of the Hamiltonian projected onto the complex,
fixed-momentum tangent space of a single-field relativistic continuous matrix
product state (RCMPS). Four auxiliary channels account for zero, ket, bra, and
joint tangent insertions. Momentum-shifted Poisson equations supply the
infinite-line boundaries, and a rank-one subtraction removes the extensive
vacuum contribution before numerical integration. The normal-ordered source
ODEs and kinetic resolvent quadrature of the refactored ground-state code lift
to these channels without a spatial lattice cutoff. We give the additional
kinetic insertion, the physical tangent metric, verification identities, and
limitations of the small-bond dense eigensolver.
\end{abstract}

\section{Scope, literature, and conventions}
The fixed-momentum tangent ansatz and its gauge quotient originate in the MPS
and cMPS excitation methods; see \cite{haegeman,draxler,vanderstraeten}.
The relativistic change of physical modes, normal ordering, and the mass-one
Hamiltonian used here follow \cite{tilloy}. Its discussion of excitation
spectra motivates the calculation but does not provide the RCMPS environments
implemented below. The four-channel connected construction and its extension
to the refactored source and resolvent engines are derived explicitly here.

We assume a translation-invariant left-canonical state with
\begin{align}
 Q&=-iK-\tfrac12R^\dagger R, & K&=K^\dagger,\\
 \LL(X)&=QX+XQ^\dagger+RXR^\dagger,
 &\LL(\rho)&=0,\quad \tr\rho=1,\quad \LL^*(\Id)=0.
\end{align}
All other transfer eigenvalues must have strictly negative real part. A
unique stationary state alone is insufficient if other peripheral modes
exist. Define
\begin{equation}
 \PP X=\rho\tr X,\qquad
 \GG=(-\LL+\PP)^{-1}(1-\PP)
       =\int_0^\infty(e^{x\LL}-\PP)\,dx.
 \label{eq:G}
\end{equation}
The matrices in the implementation are complex, column-major, and satisfy
\begin{equation}
 \operatorname{vec}\LL(X)=
 [\Id\otimes Q+\overline Q\otimes\Id+\overline R\otimes R]
 \operatorname{vec}X.
\end{equation}
The adjoint acts as
$\LL^*(A)=Q^\dagger A+AQ+R^\dagger AR$, so that
$\tr[A\LL(X)]=\tr[\LL^*(A)X]$ even when $A$ is not Hermitian.

Bare mass and target speed of light are one. The canonical physical modes
satisfy $[a(x),a^\dagger(y)]=\delta(x-y)$ and
\begin{equation}
 \phi(x)=\int J(x-y)[a(y)+a^\dagger(y)]\,dy,\qquad
 \widehat J(k)=\frac1{\sqrt{2\omega_k}},\quad
 \omega_k=\sqrt{1+k^2}.
 \label{eq:J}
\end{equation}
We use $\widehat J(k)=\int e^{-ikx}J(x)dx$. In particular,
\begin{equation}
 J(x)=\frac{K_{1/4}(|x|)}{2^{9/4}\sqrt\pi\,\Gamma(5/4)|x|^{1/4}},
 \qquad \lim_{t\to0}2|t|J(t^2)=\frac1{\sqrt\pi}.
\end{equation}
These are precisely \texttt{jsource} and \texttt{transformed\_source}.

\section{Complex tangent space and gauge fixing}
Replace one $Q,R$ pair by $V,W$ at position $x$, and integrate the resulting
state against $e^{ipx}$. We denote it by $\ket{\Phi_p(V,W)}$. Path ordering is
chosen so that auxiliary propagation acts as $X\mapsto e^{x\LL}X$. An
unambiguous realization, also fixing the sign of momentum, is the lower
triangular ket lift
\begin{equation}
 \widetilde Q_W=\begin{pmatrix}Q&0\\V&Q-ip\Id\end{pmatrix},\qquad
 \widetilde R_W=\begin{pmatrix}R&0\\W&R\end{pmatrix}.
 \label{eq:lift}
\end{equation}
The lower-left component of its ordered exponential contains one insertion
weighted by $e^{ipx}$, up to an irrelevant phase at the right endpoint.
The bra uses the analogous matrices with $U,V_U$ at the \emph{same} $p$,
followed by Hermitian conjugation in the contraction.

Pure gauge insertions are
\begin{equation}
 (V,W)=([X,Q]+ipX,[X,R]).
\end{equation}
Choose the left gauge
\begin{equation}
 V=-R^\dagger W,\qquad V_U=-R^\dagger U.
 \label{eq:gauge}
\end{equation}
For $p=0$ this also excludes the ground-state component. Define the single
and coincident insertion maps
\begin{align}
 A_W(X)&=VX+WXR^\dagger,\\
 B_U(X)&=XV_U^\dagger+RXU^\dagger,\\
 C_{UW}(X)&=WXU^\dagger.
\end{align}
The identities $\tr A_W(X)=\tr B_U(X)=0$ hold for every $X$.
Consequently, separated insertions vanish in the tangent overlap and
\begin{equation}
 \bra{\Phi_p(U)}\Phi_q(W)\rangle
 =2\pi\delta(p-q)N_{UW},\qquad N_{UW}=\tr(W\rho U^\dagger).
 \label{eq:metric}
\end{equation}
This is a complex Hilbert-space metric. Neither the factor $2\operatorname{Re}$
nor the optimizer's diagonal regularizer belongs in this eigenproblem.
The projected Hamiltonian is a sesquilinear form, not the real Hessian of
the variational energy. A real Hessian also contains two-ket/two-bra terms
and curvature terms from a chosen parameter retraction.

\section{Four channels and the extensive term}
Write a mixed lifted density as four $D\times D$ matrices
$\bm X=(X_{00},X_{10},X_{01},X_{11})$. Expanding
$\widetilde Q_W X+X\widetilde Q_U^\dagger+
 \widetilde R_W X\widetilde R_U^\dagger$ gives
\begin{align}
 \dot X_{00}&=\LL X_{00},\\
 \dot X_{10}&=(\LL-ip)X_{10}+A_W(X_{00}),\\
 \dot X_{01}&=(\LL+ip)X_{01}+B_U(X_{00}),\\
 \dot X_{11}&=\LL X_{11}+A_W(X_{01})+B_U(X_{10})+C_{UW}(X_{00}).
 \label{eq:uncentered}
\end{align}
The two shifted single-insertion environments are trace-free:
\begin{align}
 r_{10}&=(-\LL+ip+\PP)^{-1}A_W(\rho),\\
 r_{01}&=(-\LL-ip+\PP)^{-1}B_U(\rho).
 \label{eq:singles}
\end{align}
Here $\PP$ fixes the zero-momentum gauge and is harmless for every nonzero
$p$, since the right-hand sides and solutions are trace-free. In particular,
there is no numerical $1/p$ subtraction close to zero momentum.

Set
\begin{equation}
 F=A_W(r_{01})+B_U(r_{10})+W\rho U^\dagger,
 \qquad \tr F=N_{UW}.
\end{equation}
A long empty interval of length $b$ produces
$X_{11}=bN_{UW}\rho+r_{11}+o(1)$, where
\begin{equation}
 r_{11}=\GG F,\qquad
 -\LL r_{11}=F-N_{UW}\rho,\quad \tr r_{11}=0.
 \label{eq:double}
\end{equation}
The secular term is the tangent overlap proportional to the interval length.
Multiplying it by the ground-state energy density gives the unphysical
extensive part of $\bra{\Phi}H\ket{\Phi}$.

Remove it analytically. Define a connected transfer generator $\mathbb L$
by retaining the first three equations of \eqref{eq:uncentered} and replacing
the last one with
\begin{equation}
 (\mathbb L\bm X)_{11}=\LL X_{11}+A_W(X_{01})+B_U(X_{10})
                         +C_{UW}(X_{00})-N_{UW}X_{00}.
 \label{eq:connected}
\end{equation}
Then
\begin{equation}
 \bm r=(\rho,r_{10},r_{01},r_{11}),\qquad \mathbb L\bm r=0.
\end{equation}
The subtraction is the identity map $-N_{UW}X_{00}$, not the projector
$-N_{UW}\rho\tr X_{00}$. This distinction matters inside an observable source.

\section{The second boundary: the right-tail environment}
Using $\tr A=\tr B=0$, the left zero functional of $\mathbb L$ has the form
\begin{equation}
 \ell(\bm X)=\tr X_{11}+\tr(h_{UW}X_{00}),
 \quad
 -\LL^*h_{UW}=U^\dagger W-N_{UW}\Id,
 \quad \tr(\rho h_{UW})=0.
 \label{eq:left}
\end{equation}
Thus $\ell\mathbb L=0$ and $\ell(\bm r)=0$. This is the same adjoint Poisson
equation used by the refactored \texttt{boundary\_environment}, now with a
generally non-Hermitian right-hand side. Do not Hermitize $h_{UW}$ for
an off-diagonal matrix element.

For an observable supported in a finite source window, with connected
propagator $\mathbb U_O$, its excitation matrix element is
\begin{equation}
 \boxed{\quad O^{\rm exc}_{UW}(p)=\ell\bigl(\mathbb U_O\bm r\bigr).\quad}
 \label{eq:observable}
\end{equation}
In particular the identity gives zero. To see the role of the second term,
let $Y_{00}$ be the outgoing ground channel and $e_O=\tr Y_{00}$.
The tangent insertions to the right of the source contribute
\begin{equation}
 \int_0^\infty \tr\{U^\dagger W
          [e^{x\LL}Y_{00}-e_O\rho]\}\,dx
 =\tr\{U^\dagger W\GG Y_{00}\}
 =\tr(h_{UW}Y_{00}).
\end{equation}
The singles and $r_{11}$ account for insertions to the left. Hence the
source window need only resolve the smearing kernel; it need not span the
possibly much larger auxiliary correlation length.

Equivalently, insert $O$ in a box $[-B,B]$, start in the ground block only,
use the uncentered lifted transfer, and form
\begin{equation}
 O^{\rm exc}_{UW}(p)=\lim_{B\to\infty}
 \left[\bra{\Phi_{p,B}(U)}O\ket{\Phi_{p,B}(W)}
                  -2BN_{UW}\bra\Psi O\ket\Psi\right].
 \label{eq:box}
\end{equation}
Mixing gives \eqref{eq:observable} by applying
\eqref{eq:singles}--\eqref{eq:left} to the two empty tails.
Translational invariance converts the two position integrals at fixed
$O=h(0)$ into the Hamiltonian density matrix element per $2\pi\delta(0)$:
\begin{equation}
 \bra{\Phi_p(U)}(H-E_0)\ket{\Phi_q(W)}
   =2\pi\delta(p-q)H^{\rm exc}_{UW}(p).
\end{equation}
This derivation needs no exact ground-state equation. Physical interpretation
as a particle spectrum does need an adequate optimized background.

\section{Lifting the RCMPS source equations}
For a base insertion $c$ and its tangent replacements $c_W,c_U$, define
\begin{align}
 (S_L\bm X)_{00}&=cX_{00},& (S_R\bm X)_{00}&=X_{00}c^\dagger,\\
 (S_L\bm X)_{10}&=cX_{10}+c_WX_{00},&
 (S_R\bm X)_{10}&=X_{10}c^\dagger,\\
 (S_L\bm X)_{01}&=cX_{01},&
 (S_R\bm X)_{01}&=X_{01}c^\dagger+X_{00}c_U^\dagger,\\
 (S_L\bm X)_{11}&=cX_{11}+c_WX_{01},&
 (S_R\bm X)_{11}&=X_{11}c^\dagger+X_{10}c_U^\dagger.
 \label{eq:sources}
\end{align}
For a normal-ordered vertex $:e^{\beta\phi(0)}:$, take
$(c,c_W,c_U)=(R,W,U)$ and solve
\begin{equation}
 \frac{d\bm X}{dt}=2|t|\mathbb L\bm X+
       \beta\,[2|t|J(t^2)](S_L+S_R)\bm X,
 \qquad x=t|t|.
\end{equation}
Both source coefficients are $\beta$, including for imaginary $\beta$;
conjugating one of them would change the operator. A single complex vertex
is not Hermitian, but the sine-Gordon combination is.

For $:\phi^n:$ use channels $\bm X^{(0)},\ldots,\bm X^{(n)}$ with
\begin{equation}
 \dot{\bm X}^{(k)}=2|t|\mathbb L\bm X^{(k)}
       +k[2|t|J(t^2)](S_L+S_R)\bm X^{(k-1)}.
\end{equation}
Initialize only $\bm X^{(0)}=\bm r$ and contract the last channel with $\ell$.
This is exactly the refactor's \texttt{Observable}/\texttt{Edge} graph with
four tangent channels attached to each original observable channel.
The $x=t|t|$ transformation removes the integrable kernel singularity at
zero; an ODE stop at zero handles its remaining nonsmoothness.

\section{Kinetic insertion and exact convolution}
The free Hamiltonian density integrates to the sum of two bilinears, with
base insertions $c=R$ and $c=[Q,R]$. The second must be differentiated
including its spatial dependence. Taking the commutator of the lifted
matrices in \eqref{eq:lift} gives
\begin{equation}
 \boxed{c_W=[Q,W]+[V,R]-ipW,\qquad
        c_U=[Q,U]+[V_U,R]-ipU.}
 \label{eq:kinetic}
\end{equation}
For the first term, $c_W=W,c_U=U$. The sign is tied to
\eqref{eq:lift}; reversing the Fourier convention reverses it and all
momentum shifts together. The pure gauge pair gives
$c_W=[X,[Q,R]]$, which is a useful covariance check. Omitting $-ipW$
would already miss the $p^2$ term for a free-vacuum excitation.

The refactor's four-channel kinetic \emph{observable} graph, with two
$\sqrt2J$ sources, can be used directly with \eqref{eq:sources}. This is the
finite-source-window ODE backend. For an infinite-line calculation,
integrate the convolution first:
\begin{equation}
 (J*J)(x)=\frac{K_0(|x|)}{2\pi},\qquad
 K_0(x)=\int_0^\infty e^{-x\cosh u}\,du\quad(x>0).
\end{equation}
Changing variables to $\cos\theta=\operatorname{sech}u$ gives the kinetic
matrix element
\begin{equation}
 \boxed{K^{\rm exc}_{UW}[c]=\frac1\pi\int_0^{\pi/2}\!
 \ell\left[
 S_L(1-\cos\theta\,\mathbb L)^{-1}S_R\bm r+
 S_R(1-\cos\theta\,\mathbb L)^{-1}S_L\bm r
 \right]d\theta.}
 \label{eq:kinetic-resolvent}
\end{equation}
For equal tangents the two orders are complex conjugates, but for arbitrary
$U,W$ one must keep both, rather than take the real part of one order.
The existing \texttt{kinetic\_grid} supplies weights for
$(2/\pi)d\theta$, hence each order has weight one half in code.
Node doubling checks the full complex kinetic matrix element.

There is no need to form a $4D^2\times4D^2$ matrix. For
$\bm Y=(1-z\mathbb L)^{-1}\bm B$, solve consecutively
\begin{align}
 Y_{00}&=(1-z\LL)^{-1}B_{00},\\
 Y_{10}&=[1-z(\LL-ip)]^{-1}[B_{10}+zA_W(Y_{00})],\\
 Y_{01}&=[1-z(\LL+ip)]^{-1}[B_{01}+zB_U(Y_{00})],\\
 Y_{11}&=(1-z\LL)^{-1}\bigl[B_{11}+z\{A_W(Y_{01})+B_U(Y_{10})
                         +C_{UW}(Y_{00})-N_{UW}Y_{00}\}\bigr].
\end{align}
A single Schur factorization of the refactored transfer matrix supplies all
four solves at every quadrature node. Defective transfer matrices do not
require an eigenvector inverse. This backend has no smearing-window error
for the kinetic energy; potential terms still use a finite source window.

\section{Hamiltonians and the finite eigenproblem}
The interaction term lists are shared with the ground-state objectives:
\begin{align}
 H_{\phi^4}&=H_0+g\int :\phi^4:\,dx,\\
 H_{\mathrm{shG/sG}}&=H_0+\int\left[
 \frac{:e^{\beta\phi}:+:e^{-\beta\phi}:}{2\beta^2}
              -\frac12:\phi^2:\right]dx.
\end{align}
Here $\beta$ is real for shG and imaginary for sG, and additive constants
cancel. The modified-mass model uses precisely the coefficients of
\texttt{fg\_phi4\_massmod}; its energies are in that objective's rescaled
units. In particular, its kinetic prefactor is $1/m_{\rm mod}^2$, so a
speed-one diagnostic cannot be applied in those units without conversion.

Diagonalize $\rho=U_\rho\operatorname{diag}(r_a)U_\rho^\dagger$ and retain
its resolved positive support. With
\begin{equation}
 W_{ia}=\frac{\ket i\bra{u_a}}{\sqrt{r_a}},
 \qquad \tr(W_{ia}\rho W_{jb}^\dagger)=\delta_{ij}\delta_{ab},
\end{equation}
build $H_{ia,jb}=H^{\rm exc}_{W_{ia},W_{jb}}$ and solve the ordinary
Hermitian eigenproblem. Resolved support gives $D\operatorname{rank}\rho$
directions. A relative eigenvalue threshold removes unresolved or null
metric directions; the retained rank and discarded density weight are
reported. This is a restriction of the variational subspace, not an additive
metric regularization; its sensitivity must be checked for nearly singular
$\rho$.

Every matrix entry is evaluated independently. Before symmetrizing numerical
roundoff, the implementation checks
$\|H-H^\dagger\|/\max(1,\|H\|)$ and throws on an excessive defect.
Returned Ritz residuals measure the finite matrix eigenproblem only; they
are not residuals of the full quantum-field Hamiltonian. Negative levels
are retained, since clipping them would conceal a poor background or an
instability. The current code does not certify ground-state stationarity.

The rest energy is the lowest represented level at $p=0$.
It is computed by \texttt{excitation\_gap}, with the equivalent explicit-model
overload \texttt{mass\_gap(s,model)}.
The function \texttt{transfer\_gap} measures an auxiliary inverse correlation length;
the old one-argument \texttt{mass\_gap(s)} remains its compatibility alias.
These quantities are not interchangeable.
For an isolated relativistic particle of mass $M$, test
\begin{equation}
 E(p)^2-p^2=M^2
\end{equation}
and fit $E^2=a+c_{\rm fit}^2p^2$. A dispersion grid supplies energy-sorted
levels independently at each momentum. An index is not a guaranteed particle
label through a crossing. Neither a good fit nor a small eigenproblem
residual proves Lorentz invariance.

\section{Verification and practical limits}
The tests include the following independent checks.
\begin{enumerate}
\item For $D=1,R=0$, $A=B=0$, all boundary corrections vanish, and the direct
kinetic insertions give
$2(1+p^2)|\widehat J(p)|^2=\sqrt{1+p^2}$.
The infinite-line resolvent reproduces this dispersion to floating-point
precision, with rest gap one. A finite window replaces $\widehat J$ by its
truncated Fourier integral.
\item For a coherent background $R=\alpha$ and
$v=\sqrt2\operatorname{Re}\alpha$, normal ordering and displacement give
\begin{align}
 E_{\phi^4}(p)&=\omega_p+6gv^2/\omega_p,\\
 E_{\mathrm{shG/sG}}(p)&=\omega_p+
                  [\cosh(\beta v)-1]/(2\omega_p).
\end{align}
These are projected tangent energies, even when the coherent background is
not stationary; they are not the full Bogoliubov spectrum.
\item Random complex off-diagonal pairs test sesquilinearity, Hermiticity,
constant-density cancellation, the two kinetic backends, and the zero-mode
and shifted environment equations.
\item A separate reference constructs ordinary $2D\times2D$ ket and bra
matrices and their full Kronecker transfer. It evolves a finite box without
any Poisson solve or connected subtraction, then subtracts $2BN_{UW}e_O$
explicitly as in \eqref{eq:box}. Increasing the empty tails recovers the
connected monomial and kinetic-source matrix elements, at zero and nonzero
momentum. Thus agreement is not imposed by reusing the environment equations.
\item Auxiliary unitary gauge transformations leave the spectra invariant;
Ritz vectors are normalized in the physical metric. Null metric directions,
invalid controls, and unsuccessful quadrature/ODE solves are exercised.
\end{enumerate}

The present implementation is a \emph{small-bond reference}:
one dense transfer factorization costs $O(D^6)$ time and $O(D^4)$ storage;
a full tangent Hamiltonian needs $O(D^4)$ bilinear observable evaluations.
A source RHS uses $O(nD^3)$ work and $O(nD^2)$ scratch for $n$ observable
channels. The dense solver defaults to a guard at $D=16$, which is a guard,
not a promise of practical performance at that size. Increasing it explicitly
accepts the cost. The separate reverse-contraction action described below avoids the
bilinear assembly, but retains dense transfer factors; a fully matrix-free
transfer backend remains future work.

The source window, ODE tolerances, kinetic node doubling, metric support
threshold, background optimization, and bond dimension are distinct error
controls. The environment solves include the auxiliary tails exactly up to
linear-solver accuracy. Window doubling is a user convergence check, not an
automatic guarantee. Large momenta require finer ODE resolution and more
kinetic quadrature nodes. A fixed random small-$D$ state's spectrum is not
a quantitative interacting-QFT prediction.

Only excitations with the same vacuum on both spatial ends are included.
Kinks between different broken-symmetry vacua and general multiparticle
states require a larger ansatz. Thus the returned rest gap is a gap in the
represented tangent sector; it need not be the smallest mass in every
superselection sector, nor the minimum energy over all momenta.


\section{Reverse contraction: a direct Hamiltonian action}
The original bilinear calculation needs one evaluation for every pair of
basis vectors. The additional file \texttt{excitation\_action.jl} instead
returns a matrix $F_W$ such that
\begin{equation}
 H^{\rm exc}_{UW}=\tr(U^\dagger F_W).
\end{equation}
This is an analytic reverse contraction of the same four channels, not a
finite-difference Hessian. Only the ground and ket channels $x_0,x_1$ need
forward propagation. Backward covectors $a_0,a_1$ multiply the bra-only and
joint channels respectively. For an empty interval they obey, in reversed
spatial coordinate,
\begin{align}
 \dot a_1&=\LL^*a_1,\\
 \dot a_0&=(\LL^*+ip)a_0+A_W^{\rm d}(a_1),
 &A_W^{\rm d}(a)&=aV+R^\dagger aW.
\end{align}
The terminal covectors are $(0,\Id)$ on the final observable channel.
Source edges are transposed using the trace pairing, keeping their complex
coefficients unchanged. Introduce
\begin{equation}
 B^{\rm pull}(a,x)=(aR-Ra)x,\qquad W_\rho=W\rho.
\end{equation}
At every observable channel the empty-interval contribution to $F_W$ is
\begin{equation}
 B^{\rm pull}(a_0,x_0)+B^{\rm pull}(a_1,x_1)
           +a_1Wx_0-\tr(a_1x_0)W_\rho.
 \label{eq:actionbulk}
\end{equation}
For a right-source edge from $k$ to $l$, put
$P=a_0^{(l)}x_0^{(k)}+a_1^{(l)}x_1^{(k)}$. Its bra replacement contributes
$P$ for $c=R$, and
\begin{equation}
 Q^\dagger P-PQ^\dagger-RPR^\dagger+RR^\dagger P+ipP
\end{equation}
for $c=[Q,R]$, multiplied by the original edge coefficient and smearing
kernel. This expression follows by moving $U^\dagger$ to the front in
$\tr(Pc_U^\dagger)$. Left-source replacements already enter the forward ket
and backward equations.

Both boundaries can also be contracted without a chosen $U$. At the left
endpoint define
\begin{align}
 d&=\GG^{\rm d}a_1,\\
 e&=(-\LL^*-ip+\PP^*)^{-1}(1-\PP^*)
                         [a_0+A_W^{\rm d}d],
\end{align}
where $\PP^*a=\Id\tr(\rho a)$ and the superscript d denotes the trace dual.
The initial and final boundary terms in the action are
\begin{equation}
 B^{\rm pull}(e,\rho)+B^{\rm pull}(d,r_{10})+dW\rho
                        +W\GG x_{0,\rm out}.
 \label{eq:actionbound}
\end{equation}
Equations \eqref{eq:actionbulk}--\eqref{eq:actionbound}, integrated using the
refactor's existing quadrature and source transformation, give the action.
Small-system tests compare its contraction with arbitrary complex $U$
against every bilinear source type, both kinetic backends, and both signs
of complex vertex coefficients.

The kinetic resolvent is differentiated by the same reverse contractions.
Each forward solve in the two retained channels has a two-channel reverse
solve using the adjoint Schur factors. Its bra-transfer contribution is
$z$ times \eqref{eq:actionbulk}, $z=\cos\theta$; the first and second source
replacements and \eqref{eq:actionbound} complete the derivative.
The quadratic observable admits the additional infinite-line formula
\begin{equation}
 (:\phi^2:)^{\rm exc}=\frac1\pi\int_0^{\pi/2}
 \ell\left[(S_L+S_R)(1-\cos\theta\,\mathbb L)^{-1}
                         (S_L+S_R)\bm r\right]d\theta.
\end{equation}
Thus its action uses all four source orders, each with one half of the
existing kinetic-grid weight. This removes a second stiff potential ODE
in the sinh-Gordon calculation. The remaining vertex source can be evaluated
by the adaptive ODE. The D=16 quartic research run uses an independently checked
fourth-order implicit Gauss collocation rule with reused factors. Refining
the Gauss grid and checking the final Rayleigh quotient with the adaptive
ODE separate discretization error from the variational error.

For a Z$_2$-symmetric state with $ZRZ=-R$, $ZQZ=Q$ and $Z^2=\Id$,
physical parity acts on the horizontal matrix as $W\mapsto-ZWZ$.
Consequently block-diagonal $W$ gives odd physical parity, while off-diagonal
$W$ gives even parity. In either fixed-parity sector the two vertex matrix
elements at $\beta$ and $-\beta$ coincide after projection. One vertex
therefore suffices, with doubled coefficient. This saves work without
approximating the potential or mixing the parity sectors.

\section{Implicit integration used for the quartic mass calculation}
The research driver retains the continuum ansatz and discretizes only the
environment integration. With $x=t|t|$, split $t\in[-\sqrt{L},\sqrt{L}]$
into an even number $n$ of intervals, so that the origin is an interval
boundary. The two-stage Gauss coefficients are
\begin{equation}
 A=\begin{pmatrix}1/4&1/4-\sqrt{3}/6\\
                   1/4+\sqrt{3}/6&1/4\end{pmatrix},\qquad
 b_1=b_2=1/2,\qquad c_{1,2}=1/2\mp\sqrt{3}/6.
\end{equation}
For each monomial channel the vectorized equation has the form
$\dot y=M(t)y+f(t)$, where $M(t)=2|t|L_{\rm mat}$ and $f$ depends on
previous observable channels and, for the ket, the ground channels.
The stage slopes solve
\begin{equation}
 \sum_j(\delta_{ij}I-h A_{ij}M_i)k_j=M_i y+f_i,
 \qquad Y_i=y+h\sum_j A_{ij}k_j,\qquad
 y_+=y+\frac h2(k_1+k_2).
\end{equation}
The same $2D^2\times2D^2$ factorization serves every channel and tangent
vector. Reverse propagation uses its adjoint. The symplectic identity
$b_iA_{ij}+b_jA_{ji}=b_ib_j$ makes the matched Gauss stages a discrete
adjoint pair. This preserves the bilinear Hermiticity much more accurately
than the integration error itself; Hermiticity alone is therefore not a
convergence test. The research checks compare successive $n$ and the final
eigenvector's Rayleigh quotient and action against the independent adaptive
ODE implementation. $L$ is a source integration window, not a finite-volume
quantization length.

Within either parity block, diagonalize $\rho=U\Lambda U^\dagger$ and set
$B=U\Lambda^{-1/2}$. Basis matrices $W_{ij}=E_{ij}B^\dagger$ have unit
physical norm, and a Hamiltonian column is read from $F_{W_{ij}}B$.
There are $D^2/2$ such directions for balanced auxiliary parity blocks.
With fixed integration sizes, the present dense implementation has leading
work $O(nD^6)$ for factorization and full sector assembly. Its retained LU
factors and two stage matrices occupy approximately $96nD^4$ bytes in
complex double precision, excluding work arrays. These costs belong to the
present reference implementation, not to a lower bound on the tangent
method; iterative eigensolvers and different environment solvers can change
them substantially.

\paragraph{Quadrature without a factorization cache.}
The public \texttt{assembly=:action} route instead uses adaptive ODE
propagation and the same Gauss--Legendre quadrature as the refactored ground
gradient. For the quartic action, the two saved trajectories contain two
tangent channels and five monomial channels at $2q+1$ points. Their leading
storage is therefore $640qD^2$ bytes per concurrent action, not $96nD^4$.
The shared dense transfer context still retains approximately $96D^4$ bytes.
With fixed ODE step counts and quadrature orders, full parity-sector
assembly has an $O(D^5)$ potential contribution and an $O(D^6)$ contribution
from the present dense kinetic resolvents. Increased stiffness with bond
dimension can raise the step counts, so these powers provide extrapolations
rather than guaranteed runtimes.

To remove the fixed-grid integration error from an isolated lowest Ritz
pair, the research driver recomputes $Hc$ with the adaptive action and
corrects $c$ using the inverse of the grid Hamiltonian on its higher-mode
subspace. The grid matrix then serves only as a preconditioner. Each reported
refined energy is a freshly evaluated continuum Rayleigh quotient; the
original four-channel bilinear solver supplies a further independent check.
The eigenvector residual and changes under tolerance, quadrature, and window
refinement distinguish discretization error from the attainable floating-point
floor. These numerical tests do not bound the finite-$D$ physical error.

\section{Iterative lowest-eigenvalue calculation}
The white coordinates above also give an operator that can be passed directly
to a Hermitian Lanczos solver. If $\mathcal W c$ denotes conversion from
coordinates to $W$, define
\begin{equation}
 \mathsf Hc=\bigl[F_{\mathcal Wc}B\bigr]_{\rm sector}.
\end{equation}
An application needs only the reverse environment action, not a Hamiltonian
column cache. Starting from $W=\Id$ selects the odd sector in the balanced
Z$_2$ gauge and uses no previously solved excitation. This start is a trial
vector, not an assertion that the identity tangent is an eigenstate.
The forward and backward source trajectories are independent; they can be
computed concurrently, as can the kinetic and potential contributions.
For every returned vector, the research driver recomputes a fresh action to
check $\|\mathsf Hc-Ec\|$, then repeats the check with tighter integration.
This is necessary because roundoff and ODE errors can invalidate a purely
recursive Lanczos residual estimate.

A more physical zero-momentum initial vector is obtained from
$\phi(0)|\Psi\rangle$. Apart from its scalar smearing factor, annihilation
inserts $V=R$ and creation inserts $W=\Id$, so its raw pair is $(R,\Id)$.
Under the zero-momentum gauge transformation
$(V,W)\mapsto(V+[Q,X],W+[R,X])$, the left-gauge defect changes by
$-\mathcal L^\dagger X$. Consequently solve
\begin{equation}
 \mathcal L^\dagger X=R+R^\dagger,\qquad
 \operatorname{tr}(\rho X)=0,
 \qquad W_\phi=\Id+[R,X].
\end{equation}
In an exactly Z$_2$-symmetric state the right-hand side has zero stationary
expectation, so the projected Poisson solve is well-defined. The resulting
$V_\phi=R+[Q,X]$ obeys $V_\phi+R^\dagger W_\phi=0$.
Normalize its white coordinates before iteration. This construction uses
only the ground state, not an already computed excitation.

For poorly conditioned tangent metrics, the research driver also supports
full orthogonalization with the actual projected matrix $V^\dagger A$,
$A=[\mathsf Hv_1,\ldots,\mathsf Hv_k]$. With no preconditioner and before
restart this builds the same Krylov space as Hermitian Lanczos in exact
arithmetic. It measures the anti-Hermitian part before diagonalizing the
Hermitian part of the small projected matrix; it does not replace the
physical action by a presumed Hermitian recurrence. The residual uses
$Ay-EVy$, followed at termination by a fresh action on $Vy$.
Finite ODE and floating-point errors mean that these two residuals need
not agree. In particular, whitening amplifies errors along small
eigenvalues of $\rho$; retaining every positive metric direction does not
guarantee a machine-precision eigenvalue calculation.

At each kinetic quadrature node, the triangular Schur resolvent factor
is shared by its forward and adjoint solves. At $p=0$ the momentum-shifted
factor is the same matrix as the unshifted one. Reusing these matrices
reduces temporary allocations without changing the quadrature or retaining
a factorization cache over all nodes.

A second variant retains the quadratic excitation matrix $\mathsf K$ but
never assembles the interacting potential matrix. For
$\mathsf K=U\operatorname{diag}(\kappa_j)U^\dagger$, use the fixed positive
preconditioner
\begin{equation}
 P^{-1}r=U\operatorname{diag}\!\left(\frac1{\max(1,\kappa_j)}\right)U^\dagger r.
\end{equation}
The floor is the reference energy scale $m=1$, not a modification of the
physical Hamiltonian. A Davidson iteration builds orthonormal columns $V$,
stores their actions $A=\mathsf HV$, and solves the small Rayleigh--Ritz
problem $V^\dagger A y=Ey$. With $c=Vy$ and $r=Ay-Ec$, the next direction is
the twice-orthogonalized vector $-(1-VV^\dagger)P^{-1}r$. At restart the
lowest few Ritz vectors and their already computed actions are retained.
The reported energy always uses $\mathsf H$, never $P$.

The preconditioner need not use a converged kinetic quadrature: a low-order
approximation $\widetilde{\mathsf K}$ can replace $\mathsf K$ in this
construction after taking its Hermitian part and applying the same positive
spectral floor. This reduces its setup cost. The full action used for every
Ritz pair and final residual must still use the accurate kinetic resolvents
and adaptive potential environments. Thus coarsening the preconditioner
changes convergence speed, not the eigenproblem being solved. Its low-order
eigenvalues have no independent physical interpretation.

Near convergence, accumulated action errors can leave a fresh residual
larger than the stored subspace residual. A bounded residual refresh uses
$r=\mathsf Hc-Ec$ from a new, more accurate action. Set
$d\propto(1-cc^\dagger)P^{-1}r$ with $\|d\|=1$, evaluate
$s=(\mathsf H-E)d$, and choose
\begin{equation}
 \alpha=\frac{s^\dagger r}{s^\dagger s},\qquad
 c_{\rm trial}=\frac{c-\alpha d}{\|c-\alpha d\|}.
\end{equation}
This minimizes the linearized residual $\|r-\alpha s\|$ at fixed $E$.
The nonlinear, normalized eigenvector residual must be recomputed with
another fresh action; retain a correction only if that measured residual
decreases. The research check caps the number of trials and also repeats
the independent bilinear Rayleigh quotient on the retained vector.
The direction action $\mathsf Hd$ may use the original integration
tolerances as an inexact proposal, since the accepted vector and its
residual are always evaluated with the tighter settings. This avoids
oversolving high-energy trial directions whose coefficient $\alpha$ is tiny.

The first variant is matrix-free with respect to the entire excitation
Hamiltonian. The preconditioned variant is matrix-free only in the quartic
part: construction or loading of $\mathsf K$ must be counted separately.
Both still use the dense transfer context described above. Neither a
successful D=16 iteration count nor removal of the excitation matrix proves
a bond-independent Krylov dimension or scalable dense transfer solves.

\subsection{Absolute residual floors and a change of integration accuracy}

Let $U$ contain orthonormal trial vectors and let $A$ contain their computed
Hamiltonian actions. The projected matrix $T=U^\dagger A$ need not be exactly
Hermitian because these actions have finite integration accuracy. Write
\begin{equation}
 T=T_{\rm h}+S,\qquad
 T_{\rm h}=\frac{T+T^\dagger}{2},\qquad
 S=\frac{T-T^\dagger}{2}.
\end{equation}
For the normalized lowest eigenvector $q$ of $T_{\rm h}$ with eigenvalue
$\lambda$, the stored full residual is $r=Aq-\lambda Uq$. Consequently,
\begin{equation}
 U^\dagger r=S q,\qquad \|r\|\geq\|S q\|.
\end{equation}
These identities hold for the computed actions, up to orthogonality and
floating-point errors. They do not require the action routine to be exactly
linear outside this finite collection of trial vectors. In particular,
a small relative defect $\|T-T^\dagger\|/\|T\|$ does not ensure a small
absolute residual on the lowest mode. Symmetrizing the projected matrix
does not remove $S q$ from the full residual.

For the archived $D=48$, $g=2.625$ checkpoint, the 30-vector subspace had
$\|T\|_{\rm F}=2370.33$ and a relative Frobenius defect
$1.3174\times10^{-8}$, but $\|S q\|_2=1.2930\times10^{-7}$.
This contribution alone exceeded the desired $10^{-7}$ stopping threshold.
It explains why stable Ritz energies cannot certify vector convergence.
Nor does it identify which individual quadrature or environment dominates
the integration error; that would require changing those controls separately.

The continuation in \texttt{research/mass\_scan/refine\_mass.jl} therefore
retains the normalized Ritz vector and the saved positive preconditioner,
then starts a new Davidson subspace with tighter actions. The kinetic
quadrature increases from 256 to 512 nodes, the observable quadrature from
512 to 768, the spatial window from 30 to 40, and the ODE tolerances from
$(10^{-13},10^{-12})$ to $(10^{-15},10^{-14})$.
The tangent coordinates and background state are unchanged. Old approximate
actions are not mixed into the new projected matrix. The final vector must
still pass a fresh action residual and independent bilinear contractions;
tighter input tolerances by themselves are not an error certificate.
All coarse checkpoints and precision handoffs are retained for comparison.

If tight actions on normalized search directions are expensive, the residual
refresh derived above can also start from a saved, not yet converged tight
Ritz vector. This is implemented in
\texttt{research/mass\_scan/refresh\_mass.jl}. Only direction proposals use
the original integration settings; every normalized candidate uses a new
tight action and must reduce the measured residual. Final acceptance requires
the tight residual and independent bilinear checks to pass, irrespective of
whether the preceding Davidson stage reached its own stopping threshold.
The preceding subspace is preserved, and the correction attempt is bounded.

The saved action must correspond to the retained coefficients, including when
the final trial is rejected. The first research implementation accidentally
captured the outer retained-action variable inside its nested action function.
A rejected trial could therefore replace the stored action without replacing
the coefficients or residual. The final file audit detected this mismatch for
three D48 points. Their independently integrated masses were evaluated on the
correct retained vectors and were unaffected. The three research drivers now
use a distinct action binding and verify $h-Ec=r$ before saving. Fresh actions
on the affected retained vectors are exported in separate verified files; raw
files and original solver flags are preserved. All twenty D16/D32 saved results
passed the same vector--action consistency audit without modification.

The original residual target must remain distinct from an explicitly
provisional mass export. On a fixed $D=48$, $g=2.625$ mode, tightening the
ODE tolerances further to $(10^{-16},10^{-15})$ changed the action vector by
$1.5526\times10^{-9}$, whereas its real Rayleigh quotient changed by only
$1.4121\times10^{-12}$. The residual norms before and after this change were
$2.4439\times10^{-7}$ and $2.4451\times10^{-7}$. For two computed actions,
\begin{equation}
 \delta E=\operatorname{Re}(c^\dagger\delta h),\qquad
 |\delta E|\leq\|\delta h\|\quad(\|c\|=1).
\end{equation}
Small scalar changes can therefore coexist with substantially larger action
changes. Both actions here correspond to exactly the same retained coefficients.
An earlier comparison against the incorrectly saved rejected-trial action gave
$3.7942\times10^{-7}$ and was invalid. The corrected experiment does not establish
an ODE precision floor near $10^{-7}$; the bounded mode calculations simply
ended above their requested residual target. This single tolerance comparison
also does not constitute a rigorous bound on every source of operator error.
The D48 plot exports explicitly preserve any failed $10^{-7}$ residual flag.
Their provisional inclusion requires residual at most $10^{-6}$, unit
normalization and parity checks, and independent scalar contractions agreeing
within $10^{-7}$. These are qualified mass estimates, not a declaration that
the stricter eigensolver threshold was reached. Ground-state optimization
bias and finite-bond effects remain separate from these numerical checks.

\subsection{Library implementation}
The package function \texttt{excitation\_eigsolve} implements the procedure of this
section; the research drivers are kept as records. \texttt{TangentCoordinates}
selects the sector ($W\mapsto-ZWZ$) and forms the white coordinates. With
\texttt{gauge=:density} it rotates each parity block into the eigenbasis of $\rho$
(an exact unitary gauge change of $(R,K,\rho)$) and uses the Cholesky factor of the
stored density, so that $B^\dagger\rho B=\Id$ holds for the $\rho$ entering the
environments. Coordinates of an arbitrary $W$ are its metric projection
$c=[W\rho B]_{\rm sector}$, and an action is read as $[F_WB]_{\rm sector}$.

The operator applies the reverse-contracted action with a fixed kinetic grid, so its
kinetic part is exactly linear; the adaptive potential ODEs are linear to their
integration tolerance. The field-insertion start vector is defined at any momentum by
\begin{equation}
 (\mathcal L^\dagger+ip)X=R+R^\dagger,\qquad W_\phi=\Id+[R,X],
\end{equation}
which reduces to the zero-momentum construction above. The shift follows from the
momentum-$p$ gauge transformation $(V,W)\mapsto(V+[Q,X]-ipX,W+[R,X])$, whose
left-gauge defect is $-(\mathcal L^\dagger+ip)X$; it is only a start vector. At $p\neq0$ the code solves the trace-projected equation, which is sufficient
because $V=-R^\dagger W$ is imposed on every tangent. One Davidson routine covers
both research solvers. Fresh restarts retain the lowest Ritz vectors and recompute
their actions; thick restarts recombine the stored action columns. Both use the same
stall rule and require the projected matrix to be Hermitian to $10^{-7}$.

The momentum-independent Schur factor and pinned Poisson factor are shared between
momenta and between integration settings; only the two $\pm ip$ factors are rebuilt.
Summing the kinetic boundary covectors of all terms before the two dual Poisson
solves is exact, because the boundary pullback is linear. Validation repeats the
research checks: fresh standard and refined actions, and two bilinear energies.

\section{Near-critical signed energies and a direct stability check}

For an approximate symmetric vacuum, the lowest odd tangent energy can be
negative even after the numerical eigenproblem is converged. Such a value
is not a positive particle mass, and taking its absolute value has no
justification. A small numerical residual controls the projected eigenproblem,
not the accuracy or unrestricted stability of the variational background.

A separate energy-density test probes the background along the computed
horizontal mode. For $W_\theta=e^{i\theta}W$, define the Hermitian matrix
\begin{equation}
 \delta K_\theta=\frac{i}{2}(W_\theta^\dagger R-R^\dagger W_\theta),\qquad
 R_\epsilon=R+\epsilon W_\theta,\quad
 K_\epsilon=K+\epsilon\delta K_\theta.
\end{equation}
Reconstructing $Q_\epsilon=-iK_\epsilon-R_\epsilon^\dagger R_\epsilon/2$
preserves left normalization exactly and gives
$\partial_\epsilon Q_\epsilon|_0=-R^\dagger W_\theta=V_\theta$.
The stationary density matrix must be recomputed at every displacement.
For an exactly symmetric background and an odd physical tangent, field
parity maps $\epsilon$ to $-\epsilon$; the even Hamiltonian therefore has
$e(\epsilon,\theta)=e(-\epsilon,\theta)$. The finite-displacement quantity
\begin{equation}
 \kappa_\theta(\epsilon)=
 \frac{e(\epsilon,\theta)-e(0)}{\epsilon^2}
\end{equation}
tests energy lowering when its sign persists as $\epsilon$ decreases and
under independent integration checks. This is an additional variational
diagnostic, not a replacement for the tangent Hamiltonian: the energy-density
Hessian can contain anomalous pair terms and background-gradient contributions.
An energy-lowering displacement outside the imposed symmetric parameter
manifold establishes an instability of that saved state, not a precise
infinite-bond critical coupling.

\begin{thebibliography}{9}
\bibitem{haegeman} J. Haegeman, J. I. Cirac, T. J. Osborne, and F. Verstraete,
\emph{Calculus of continuous matrix product states}, Phys. Rev. B 88, 085118
(2013), \href{https://arxiv.org/abs/1211.3935}{arXiv:1211.3935}.
\bibitem{draxler} D. Draxler, J. Haegeman, T. J. Osborne, V. Stojevic,
L. Vanderstraeten, and F. Verstraete,
\emph{Particles, holes and solitons: a matrix product state approach},
Phys. Rev. Lett. 111, 020402 (2013),
\href{https://arxiv.org/abs/1212.1114}{arXiv:1212.1114}.
\bibitem{vanderstraeten} L. Vanderstraeten, J. Haegeman, and F. Verstraete,
\emph{Tangent-space methods for uniform matrix product states},
SciPost Phys. Lect. Notes 7 (2019),
\href{https://scipost.org/10.21468/SciPostPhysLectNotes.7}{doi:10.21468/SciPostPhysLectNotes.7}.
\bibitem{tilloy} A. Tilloy,
\emph{Relativistic continuous matrix product states for quantum fields without cutoff},
Phys. Rev. D 104, 096007 (2021),
\href{https://arxiv.org/abs/2102.07741}{arXiv:2102.07741}.
\end{thebibliography}
\end{document}
